\begin{figure}[htbp]
\centering
\begin{tikzpicture}[x=1cm,y=1cm,>=Stealth,
 every node/.style={font=\small},
 axis/.style={draw=black!45,line width=.5pt},
 vector/.style={->,draw=blue!55!black,line width=.9pt}]
 \coordinate (p) at (0,0);
 \coordinate (y) at (2.9,.65);
 \fill (p) circle (1.5pt) node[below left] {$p$};
 \fill (y) circle (1.5pt) node[above right] {$y$};
 \draw[vector] (p) -- (y) node[midway,above,sloped] {$x=y-p$};
 \node[align=center] at (1.45,-1.05) {center and target\\in physical space};

 \begin{scope}[shift={(7.9,0)}]
  \fill[red!18] (35:1.4) arc (35:-35:1.4) -- cycle;
  \fill[red!18] (145:1.4) arc (145:215:1.4) -- cycle;
  \draw[axis] (0,0) circle (1.4);
  \draw[axis,->] (-1.8,0) -- (1.95,0) node[right] {$x/|x|$};
  \draw[dashed,black!55] (35:1.4) -- (-35:1.4);
  \draw[dashed,black!55] (145:1.4) -- (215:1.4);
  \draw[vector] (0,0) -- (20:1.4) node[above right] {$u$};
  \node[align=center] at (0,-2.05)
    {cutting normals $C(x)$\\$|u\cdot x|>1+3\eta$};
  \node at (0,1.8) {normal sphere};
 \end{scope}
\end{tikzpicture}
\caption{Plane-section schematic of a residual and its cutting cap.
The shaded antipodal caps lie beyond the two threshold hyperplanes
$u\cdot x=\pm(1+3\eta)$. A sampled normal in $C(x)$ excludes the target
$y$ from $p+K^+$. The drawing does not depict high-dimensional cap
measure. For a block of residuals, one normal in their common subset
cuts them all; assigning distinct normals to all blocks is a sufficient
certificate of a hole.}
\label{fig:residual-cap}
\end{figure}
